\begin{proof}[Proof of \cref{obj:lem:full-record-copula-component-bound}]
\begin{enumerate}
\item
We first construct the common augmentation and the three label kernels. For \(J\in\{K,M\}\), define the cell index, including its value at the right endpoint, by
\[
\operatorname{cell}_J(x)=\min\{J-1,\lfloor Jx\rfloor\},
\qquad x\in[0,1].
\]
Here \(K>0\), since \(16M\le K\) and \(M\ge2\). Define the boundary-node set and the disclosure map by
\[
\mathsf N_{\partial}
 =\{j\in\{0,\ldots,K\}:jM\bmod K=0\},
\qquad
\mathsf D(\boldsymbol\lambda,\boldsymbol\eta)_j
 =
 \begin{cases}
 (\lambda_j,\eta_j),&j\in\mathsf N_{\partial},\\
 \bot,&j\notin\mathsf N_{\partial}.
 \end{cases}
\]
The symbol \(\bot\) denotes an undisclosed pair. Thus the space of possible disclosures is
\[
\mathsf\Delta_K
 =\bigl(\{\bot\}\cup\{-1,1\}^2\bigr)^{\{0,\ldots,K\}},
\]
and the augmentation space and a generic augmentation are
\[
\mathsf\Omega_{\mathrm{aug}}
 =[0,1]^n\times\{0,1\}^n\times\mathsf\Delta_K,
\qquad
\omega=(\mathbf X,\mathbf B,\delta).
\]

Connect observations \(i\) and \(i'\) when
\[
\begin{aligned}
&\operatorname{cell}_M(X_i)=\operatorname{cell}_M(X_{i'}),\\
&\operatorname{cell}_K(X_i)=\operatorname{cell}_K(X_{i'})
\quad\text{or}\quad
\left[
\begin{array}{c}
|\operatorname{cell}_K(X_i)-\operatorname{cell}_K(X_{i'})|=1,\\
\max\{\operatorname{cell}_K(X_i),\operatorname{cell}_K(X_{i'})\}
 \notin\mathsf N_{\partial}.
\end{array}
\right]
\end{aligned}
\]
Let \(\mathscr C(\omega)\) be the collection of connected components of this graph, with connectivity including paths of length zero. It partitions \(\{1,\ldots,n\}\). For every subset \(\mathsf c\subseteq\{1,\ldots,n\}\), set
\[
\mathsf n_{\mathsf c}=|\mathsf c|,
\qquad
\mathsf k_{\mathsf c}=\sum_{i\in\mathsf c}B_i,
\qquad
\mathsf j(\mathsf c)
 =\max\left(
 \left\{\left\lfloor\frac{\operatorname{cell}_M(X_i)}2\right\rfloor:
 i\in\mathsf c\right\}\cup\{0\}\right).
\]
The last convention defines \(\mathsf j(\varnothing)=0\).

We represent the binary labels by their signs:
\[
\mathbf z=(z_i^{\mathrm t},z_i^{\mathrm y})_{i=1}^n
 \in\mathsf Z_n:=\bigl(\{-1,1\}^2\bigr)^n,
\qquad
\mathsf F_n
 =4^{-n}\sum_{\mathbf z\in\mathsf Z_n}
       \operatorname{Dirac}_{\mathbf z}.
\]
The treatment label is \((1+z_i^{\mathrm t})/2\); the outcome-sign label is retained for both values of \(B_i\).

For \(\nu\in\{0,1\}\), a coarse-sign vector
\(\sigma\in\{-1,1\}^{M/2}\), and a disclosure \(\delta\), define the conditional coefficient weights by
\[
\mathsf w_{\nu,\sigma,\delta}
 (\boldsymbol\lambda,\boldsymbol\eta)
 =
 \prod_{j=0}^{K}
 \begin{cases}
 \mathbf1\{(\lambda_j,\eta_j)=\delta_j\},&\delta_j\ne\bot,\\[2pt]
 \displaystyle
 \frac{1+\nu\kappa g_\sigma(j/K)\lambda_j\eta_j}{4},
 &\delta_j=\bot,
 \end{cases}
\qquad \kappa=\frac1{16}.
\]
For every real function \(\Psi\) on the finite coefficient space, write
\[
\mathbb E_{\nu,\sigma,\delta}[\Psi]
 =
 \sum_{\boldsymbol\lambda,\boldsymbol\eta\in\{-1,1\}^{K+1}}
 \mathsf w_{\nu,\sigma,\delta}
 (\boldsymbol\lambda,\boldsymbol\eta)\,
 \Psi(\boldsymbol\lambda,\boldsymbol\eta).
\]
Each node weight is nonnegative and sums to one: for an undisclosed node this follows from
\[
|\kappa g_\sigma(j/K)|\le\frac1{16},
\qquad
\sum_{\lambda_j,\eta_j\in\{-1,1\}}\lambda_j\eta_j=0.
\]
Consequently,
\[
\mathsf w_{\nu,\sigma,\delta}\ge0,
\qquad
\sum_{\boldsymbol\lambda,\boldsymbol\eta}
 \mathsf w_{\nu,\sigma,\delta}
 (\boldsymbol\lambda,\boldsymbol\eta)=1.
\]

Use the coordinate fields of \cref{obj:def:copula-frame-priors}, and define the record-label density by
\[
\mathsf r_{\nu,i}
 =
 1+z_i^{\mathrm t}\xi(X_i)
 +B_i\left\{
 z_i^{\mathrm y}\upsilon(X_i)
 +z_i^{\mathrm t}z_i^{\mathrm y}\zeta(X_i)
 \right\}.
\]
Its dependence on the coefficient draw and \(\sigma\) is through those coordinate fields. On each fine cell, the two active frame coordinates satisfy
\[
\sum_{j=0}^{K}f_j(x)^2=1,
\qquad
\sum_{j=0}^{K}|f_j(x)|\le\sqrt2\le2.
\]
The same identities hold at cell endpoints. They imply
\[
|\xi(x)|\le\frac18,
\qquad
|\upsilon(x)|\le\frac18,
\qquad
|\widetilde g_\sigma(x)|\le1,
\qquad
|t_\nu(x)|
 \le\frac{au\kappa}{1-a^2}\le\frac1{128}.
\]
Indeed, \(a^2\le1/256\), and hence \(1-a^2\ge255/256>0\). Moreover,
\[
|1-\xi(x)^2|\le1,
\qquad
|\zeta(x)|
 \le\frac1{64}+\frac1{128}\le\frac18.
\]
Thus, for every augmentation, coefficient draw, and label vector,
\[
\frac12\le\mathsf r_{\nu,i}\le2,
\qquad
\frac14\sum_{z_i^{\mathrm t},z_i^{\mathrm y}\in\{-1,1\}}
 \mathsf r_{\nu,i}=1.
\]

For \(s_{\mathrm{coarse}}\in\{-1,1\}\), let
\(\sigma^{s_{\mathrm{coarse}}}\) denote the constant coarse-sign vector:
\[
\sigma^{s_{\mathrm{coarse}}}_j=s_{\mathrm{coarse}}
\quad(1\le j\le M/2).
\]
Define, for every subset \(\mathsf c\),
\[
\mathsf d_{\mathsf c,\nu,s_{\mathrm{coarse}}}
 (\omega,\mathbf z)
 =
 \mathbb E_{\nu,\sigma^{s_{\mathrm{coarse}}},\delta}
 \left[\prod_{i\in\mathsf c}\mathsf r_{\nu,i}\right],
\]
and abbreviate
\[
\mathsf d_{\mathsf c,0}
 =\mathsf d_{\mathsf c,0,-1},
\qquad
\mathsf d_{\mathsf c,+}
 =\mathsf d_{\mathsf c,1,1},
\qquad
\mathsf d_{\mathsf c,-}
 =\mathsf d_{\mathsf c,1,-1},
\qquad
\mathsf d_{\mathsf c,\mathrm{av}}
 =\frac{\mathsf d_{\mathsf c,+}+\mathsf d_{\mathsf c,-}}2.
\]
The null density is independent of the coarse signs. Empty products equal one, so all these component densities equal one when \(\mathsf c=\varnothing\). Finite averaging and label normalization give
\[
2^{-\mathsf n_{\mathsf c}}
 \le\mathsf d_{\mathsf c,\nu,s_{\mathrm{coarse}}}
 \le2^{\mathsf n_{\mathsf c}},
\qquad
\int\mathsf d_{\mathsf c,\nu,s_{\mathrm{coarse}}}\,d\mathsf F_n=1,
\]
and the same bounds and normalization hold for
\(\mathsf d_{\mathsf c,\mathrm{av}}\).

Now define the three full-label densities and kernels by
\[
\begin{aligned}
\mathsf f_0(\omega,\mathbf z)
 &=\prod_{\mathsf c\in\mathscr C(\omega)}
       \mathsf d_{\mathsf c,0}(\omega,\mathbf z),\\
\mathsf f_1(\omega,\mathbf z)
 &=2^{-M/2}\sum_{\sigma\in\{-1,1\}^{M/2}}
   \mathbb E_{1,\sigma,\delta}
       \left[\prod_{i=1}^n\mathsf r_{1,i}\right],\\
\mathsf f_{\mathrm{av}}(\omega,\mathbf z)
 &=\prod_{\mathsf c\in\mathscr C(\omega)}
       \mathsf d_{\mathsf c,\mathrm{av}}(\omega,\mathbf z),\\
P_\nu(\omega,d\mathbf z)
 &=\mathsf f_\nu(\omega,\mathbf z)\,\mathsf F_n(d\mathbf z),
 \qquad \nu\in\{0,1\},\\
Q(\omega,d\mathbf z)
 &=\mathsf f_{\mathrm{av}}(\omega,\mathbf z)\,
       \mathsf F_n(d\mathbf z).
\end{aligned}
\]
These definitions apply to every augmentation, including disclosures outside the support of the augmentation law. Component densities depend only on labels in their component. Since the components partition the observations, their normalized products are normalized. The full coefficient average defining \(\mathsf f_1\) is normalized by the record-label normalization above. Therefore \(P_0(\omega),P_1(\omega),Q(\omega)\) are probability laws for every \(\omega\).

Let the disclosure law be the pushforward of independent fair coefficient pairs:
\[
\mu_{\partial}
 =
 4^{-(K+1)}
 \sum_{\boldsymbol\lambda,\boldsymbol\eta\in\{-1,1\}^{K+1}}
 \operatorname{Dirac}_{\mathsf D(\boldsymbol\lambda,\boldsymbol\eta)},
\qquad
\mu
 =\lambda^{\otimes n}
   \otimes\operatorname{Bernoulli}(\varepsilon)^{\otimes n}
   \otimes\mu_{\partial}.
\]
It is a probability law with the required covariate and flag marginals, and
\[
\delta_j\ne\bot\ \Longleftrightarrow\ j\in\mathsf N_{\partial}
\qquad\text{for }\mu\text{-almost every }\omega.
\]
The kernels are measurable: the graph is determined by finitely many measurable cell assignments, and all coefficient and label averages are finite. Their normalization gives
\[
(\mu\otimes P_\nu)\circ\operatorname{pr}_{\mathrm{aug}}^{-1}
 =\mu,\qquad \nu\in\{0,1\},
\qquad
(\mu\otimes Q)(\mathsf\Omega_{\mathrm{aug}}\times\mathsf Z_n)=1.
\]
% lean: CausalSmith.Stat.FinitepHomogeneityDensegamma.copula_common_augmentation_certificate
% lean: CausalSmith.Stat.FinitepHomogeneityDensegamma.intermediateConditionalLaw_isProbabilityMeasure

\item\label{proof:full-record-copula-component-bound:conditional-comparisons}
We establish the conditional comparisons. Write
\[
K=2^{k_{\mathrm f}},\qquad M=2^{k_{\mathrm c}},
\qquad k_{\mathrm f},k_{\mathrm c}\in\mathbb N.
\]
Since \(M\le K\) and \(M\ge2\),
\[
k_{\mathrm c}\le k_{\mathrm f},
\qquad k_{\mathrm c}\ge1,
\qquad M\mid K,\qquad 2\mid M.
\]
In particular,
\[
\mathsf N_{\partial}
 =\{jK/M:0\le j\le M\}.
\]

Fix an augmentation in the support of \(\mu\). Its disclosed coefficients are fixed. The undisclosed coefficient nodes used by a component are
\[
\mathsf N_{\mathsf c}
 =
 \{j\notin\mathsf N_{\partial}:
       f_j(X_i)\ne0\text{ for some }i\in\mathsf c\}.
\]
These node sets are disjoint for distinct components. Indeed, two observations sharing an undisclosed active node lie in the same fine interval or in neighboring fine intervals joined through that node, and therefore are connected. Each component lies in one coarse cell. Its coefficient coupling and its correction consequently depend on the single sign indexed by \(\mathsf j(\mathsf c)\). Conditional coefficient independence then gives
\[
\mathbb E_{\nu,\sigma,\delta}
 \left[\prod_{i=1}^n\mathsf r_{\nu,i}\right]
 =
 \prod_{\mathsf c\in\mathscr C(\omega)}
 \mathsf d_{\mathsf c,\nu,\sigma_{\mathsf j(\mathsf c)+1}}.
\]
Define the occupied coarse-pair indices and their component groups by
\[
\mathscr J(\omega)
 =\{\mathsf j(\mathsf c):\mathsf c\in\mathscr C(\omega)\},
\qquad
\mathscr C_j(\omega)
 =\{\mathsf c\in\mathscr C(\omega):\mathsf j(\mathsf c)=j\}.
\]
The index \(\mathsf j(\mathsf c)\) belongs to \(\{0,\ldots,M/2-1\}\). Averaging the independent fair coarse signs yields
\[
\mathsf f_1
 =
 \prod_{j\in\mathscr J(\omega)}
 \frac12\left\{
   \prod_{\mathsf c\in\mathscr C_j(\omega)}
      \mathsf d_{\mathsf c,+}
   +
   \prod_{\mathsf c\in\mathscr C_j(\omega)}
      \mathsf d_{\mathsf c,-}
 \right\}.
\]

For every augmentation and every subset \(\mathsf c\), define
\[
\Delta_{\mathsf c,\mathrm e}
 =\mathsf d_{\mathsf c,\mathrm{av}}-\mathsf d_{\mathsf c,0},
\qquad
\Delta_{\mathsf c,\mathrm o}
 =\frac{\mathsf d_{\mathsf c,+}-\mathsf d_{\mathsf c,-}}2,
\]
\[
\mathfrak e_{\mathsf c}
 =\int\frac{\Delta_{\mathsf c,\mathrm e}^2}
                 {\mathsf d_{\mathsf c,0}}\,d\mathsf F_n,
\qquad
\mathfrak o_{\mathsf c}
 =\int\frac{\Delta_{\mathsf c,\mathrm o}^2}
                 {\mathsf d_{\mathsf c,\mathrm{av}}}\,d\mathsf F_n.
\]
The denominators are strictly positive. Component normalization gives
\[
\int\Delta_{\mathsf c,\mathrm e}\,d\mathsf F_n
 =\int\Delta_{\mathsf c,\mathrm o}\,d\mathsf F_n=0,
\qquad
\mathfrak e_{\mathsf c}\ge0,\qquad
\mathfrak o_{\mathsf c}\ge0.
\]
Choose the conditional budgets
\[
\mathcal A(\omega)
 =\sum_{\mathsf c\in\mathscr C(\omega)}
       \mathfrak e_{\mathsf c},
\qquad
\mathcal B(\omega)
 =\frac12
   \sum_{\mathsf c,\mathsf c'\in\mathscr C(\omega)}
   \mathbf1\{\mathsf c\ne\mathsf c',\
             \mathsf j(\mathsf c)=\mathsf j(\mathsf c')\}
       \mathfrak o_{\mathsf c}\mathfrak o_{\mathsf c'}.
\]

For the first comparison, cancellation of the even discrepancy gives
\[
\int
 \frac{\mathsf d_{\mathsf c,\mathrm{av}}^2}
      {\mathsf d_{\mathsf c,0}}\,d\mathsf F_n
 =1+\mathfrak e_{\mathsf c}.
\]
Integration over disjoint label blocks therefore gives
\[
1+\chi^2(Q(\omega)\Vert P_0(\omega))
 =\prod_{\mathsf c\in\mathscr C(\omega)}
      (1+\mathfrak e_{\mathsf c})
 \le\exp\!\left(
       \sum_{\mathsf c\in\mathscr C(\omega)}
          \mathfrak e_{\mathsf c}\right)
 =\exp(\mathcal A(\omega)).
\]

For the second comparison, define
\[
\vartheta_{\mathsf c}
 =\frac{\Delta_{\mathsf c,\mathrm o}}
        {\mathsf d_{\mathsf c,\mathrm{av}}}.
\]
Under the component law with density
\(\mathsf d_{\mathsf c,\mathrm{av}}\),
\[
\int\vartheta_{\mathsf c}\,
       \mathsf d_{\mathsf c,\mathrm{av}}\,d\mathsf F_n=0,
\qquad
\int\vartheta_{\mathsf c}^{\,2}\,
       \mathsf d_{\mathsf c,\mathrm{av}}\,d\mathsf F_n
 =\mathfrak o_{\mathsf c}.
\]
For two signs \(s_{\mathrm{coarse}},s'_{\mathrm{coarse}}\in\{-1,1\}\), it follows that
\[
\int
 (1+s_{\mathrm{coarse}}\vartheta_{\mathsf c})
 (1+s'_{\mathrm{coarse}}\vartheta_{\mathsf c})
 \mathsf d_{\mathsf c,\mathrm{av}}\,d\mathsf F_n
 =
 1+s_{\mathrm{coarse}}s'_{\mathrm{coarse}}
       \mathfrak o_{\mathsf c}.
\]
Using two independent copies of each fair coarse sign when squaring the likelihood ratio gives the exact identity
\[
1+\chi^2(P_1(\omega)\Vert Q(\omega))
 =
 \prod_{j\in\mathscr J(\omega)}
 \frac12\left\{
 \prod_{\mathsf c\in\mathscr C_j(\omega)}
      (1+\mathfrak o_{\mathsf c})
 +
 \prod_{\mathsf c\in\mathscr C_j(\omega)}
      (1-\mathfrak o_{\mathsf c})
 \right\}.
\]

Here is the exponential bound for each factor. For any finite collection
\(\mathscr U\) of components, define
\[
\begin{aligned}
\mathsf H_{\mathrm e}(\mathscr U)
 &=\frac12\left\{
   \prod_{\mathsf c\in\mathscr U}(1+\mathfrak o_{\mathsf c})
   +\prod_{\mathsf c\in\mathscr U}(1-\mathfrak o_{\mathsf c})
   \right\},\\
\mathsf H_{\mathrm o}(\mathscr U)
 &=\frac12\left\{
   \prod_{\mathsf c\in\mathscr U}(1+\mathfrak o_{\mathsf c})
   -\prod_{\mathsf c\in\mathscr U}(1-\mathfrak o_{\mathsf c})
   \right\},\\
\mathsf D_{\mathrm o}(\mathscr U)
 &=\sum_{\mathsf c\in\mathscr U}\mathfrak o_{\mathsf c},\\
\mathsf T_{\mathrm o}(\mathscr U)
 &=\frac12\sum_{\substack{\mathsf c,\mathsf c'\in\mathscr U\\
                         \mathsf c\ne\mathsf c'}}
      \mathfrak o_{\mathsf c}\mathfrak o_{\mathsf c'}.
\end{aligned}
\]
For the empty collection,
\[
\mathsf H_{\mathrm e}=1,\qquad
\mathsf H_{\mathrm o}=\mathsf D_{\mathrm o}
 =\mathsf T_{\mathrm o}=0.
\]
Upon adjoining \(\mathsf c_*\notin\mathscr U\), the recurrences are
\[
\begin{aligned}
\mathsf H_{\mathrm e}(\mathscr U\cup\{\mathsf c_*\})
 &=\mathsf H_{\mathrm e}(\mathscr U)
   +\mathfrak o_{\mathsf c_*}\mathsf H_{\mathrm o}(\mathscr U),\\
\mathsf H_{\mathrm o}(\mathscr U\cup\{\mathsf c_*\})
 &=\mathsf H_{\mathrm o}(\mathscr U)
   +\mathfrak o_{\mathsf c_*}\mathsf H_{\mathrm e}(\mathscr U),\\
\mathsf T_{\mathrm o}(\mathscr U\cup\{\mathsf c_*\})
 &=\mathsf T_{\mathrm o}(\mathscr U)
   +\mathfrak o_{\mathsf c_*}\mathsf D_{\mathrm o}(\mathscr U).
\end{aligned}
\]
These recurrences preserve
\[
\mathsf H_{\mathrm e}\ge0,\qquad
\mathsf H_{\mathrm o}\ge0,\qquad
\mathsf H_{\mathrm o}
 \le\mathsf D_{\mathrm o}\mathsf H_{\mathrm e}.
\]
For the last inequality, the old bound gives
\[
\mathsf H_{\mathrm o}
 +\mathfrak o_{\mathsf c_*}\mathsf H_{\mathrm e}
 \le
 (\mathsf D_{\mathrm o}+\mathfrak o_{\mathsf c_*})
 \mathsf H_{\mathrm e}
 \le
 (\mathsf D_{\mathrm o}+\mathfrak o_{\mathsf c_*})
 (\mathsf H_{\mathrm e}
   +\mathfrak o_{\mathsf c_*}\mathsf H_{\mathrm o}).
\]
Induction, using \(1+x\le e^x\), now gives
\[
\begin{aligned}
\mathsf H_{\mathrm e}(\mathscr U\cup\{\mathsf c_*\})
 &\le
 \mathsf H_{\mathrm e}(\mathscr U)
 \bigl(1+\mathfrak o_{\mathsf c_*}\mathsf D_{\mathrm o}(\mathscr U)\bigr)\\
 &\le
 \exp\!\left(
 \mathsf T_{\mathrm o}(\mathscr U)
 +\mathfrak o_{\mathsf c_*}\mathsf D_{\mathrm o}(\mathscr U)\right).
\end{aligned}
\]
Applying this to each \(\mathscr C_j(\omega)\) proves
\[
1+\chi^2(P_1(\omega)\Vert Q(\omega))
 \le\exp(\mathcal B(\omega)).
\]
Finally, \(\mathsf f_0\) and \(\mathsf f_{\mathrm{av}}\) are strictly positive on the finite label alphabet. Hence
\[
Q(\omega)\ll P_0(\omega),
\qquad
P_1(\omega)\ll Q(\omega).
\]
All these comparisons are conditional on the full triple \(\omega\).
% lean: CausalSmith.Stat.FinitepHomogeneityDensegamma.copula_conditional_comparison_certificate

\item
We derive the component estimates used for occupancy counting. In this step, the finite density expressions are considered as algebraic functions of
\[
a'\in(-1,1),\qquad u'\in\mathbb R,
\]
with the augmentation and coefficient weights fixed. Quantitative derivative bounds will be used on
\([0,1/16]^2\). Write \(\Delta_{\mathsf c,\mathrm e}(a',u')\) and
\(\Delta_{\mathsf c,\mathrm o}(a',u')\) for the resulting discrepancies.

The entire propensity-coefficient marginal and the entire outcome-coefficient marginal are common to both branches, with disclosed coordinates pinned. Therefore
\[
\Delta_{\mathsf c,\mathrm e}(a',0)
 =\Delta_{\mathsf c,\mathrm o}(a',0)
 =\Delta_{\mathsf c,\mathrm e}(0,u')
 =\Delta_{\mathsf c,\mathrm o}(0,u')=0.
\]
We also need the two first-order even cancellations.

To see them, the conditional coefficient means at an undisclosed node are
\[
\mathbb E_{\nu,\sigma,\delta}
 [\eta_j\mid\boldsymbol\lambda]
 =\nu\kappa g_\sigma(j/K)\lambda_j,
\qquad
\mathbb E_{\nu,\sigma,\delta}
 [\lambda_j\mid\boldsymbol\eta]
 =\nu\kappa g_\sigma(j/K)\eta_j.
\]
At a disclosed node, these means are the corresponding pinned values. For constant coarse sign, the undisclosed means change sign when the coarse sign changes.

If \(\mathsf c\) has exactly one marked observation \(i\), its record product equals
\[
\begin{aligned}
\prod_{i'\in\mathsf c}\mathsf r_{\nu,i'}
={}&
\prod_{i'\in\mathsf c\setminus\{i\}}
 (1+z_{i'}^{\mathrm t}\xi(X_{i'}))\\
&\times
\left[
 (1+z_i^{\mathrm t}\xi(X_i))
 (1+z_i^{\mathrm y}\upsilon(X_i))
 +z_i^{\mathrm t}z_i^{\mathrm y}
   t_\nu(X_i)(1-\xi(X_i)^2)
\right].
\end{aligned}
\]
Conditioning on the propensity coefficients, its outcome-field term is affine in the outcome coefficients. The branch-dependent conditional means of the undisclosed outcome coefficients, and the correction \(t_1\), are linear in the constant coarse sign. Averaging its two values cancels these changes. Disclosed-coordinate terms are common. Thus
\[
\mathsf d_{\mathsf c,+}+\mathsf d_{\mathsf c,-}
 =2\mathsf d_{\mathsf c,0},
\qquad
\Delta_{\mathsf c,\mathrm e}=0
\quad\text{if }\mathsf k_{\mathsf c}\le1.
\]
For zero marks, every record factor uses only the propensity coefficients, so
\[
\Delta_{\mathsf c,\mathrm o}=0
\quad\text{if }\mathsf k_{\mathsf c}=0.
\]

At \(a'=0\), all undifferentiated record factors depend only on the outcome coefficients. Differentiating one factor produces a term linear in the propensity coefficients and a correction term linear in the coarse sign. Conditioning on the outcome coefficients and averaging the two constant coarse signs, using the conditional means above, therefore gives
\[
\partial_{a'}\Delta_{\mathsf c,\mathrm e}(0,u')=0.
\]
For the outcome derivative, define the modified augmentation
\[
\omega^{(i)}
 =\bigl(\mathbf X,(B_{i'}\mathbf1\{i'=i\})_{i'=1}^n,\delta\bigr).
\]
Each record factor is affine in \(u'\), so the product rule at \(u'=0\) gives
\[
\partial_{u'}\Delta_{\mathsf c,\mathrm e}(a',0;\omega)
 =
 \sum_{i\in\mathsf c}
 \left\{
 \Delta_{\mathsf c,\mathrm e}(a',1;\omega^{(i)})
 -\Delta_{\mathsf c,\mathrm e}(a',0;\omega)
 \right\}=0.
\]
Here evaluation at unit outcome amplitude extracts the coefficient of the affine marked factor; the algebraic functions are defined there. Each modified augmentation has at most one mark, so the preceding cancellation applies.

For the derivative bounds, define the coefficient fields and the rational correction factor by
\[
\mathsf F_\lambda(x)=\sum_{j=0}^{K}f_j(x)\lambda_j,
\qquad
\mathsf F_\eta(x)=\sum_{j=0}^{K}f_j(x)\eta_j,
\qquad
\mathsf c_x(a')
 =\frac{a'(1-(a')^2\mathsf F_\lambda(x)^2)}
        {1-(a')^2}.
\]
The frame identities give
\[
|\mathsf F_\lambda|,|\mathsf F_\eta|\le\sqrt2\le2,
\qquad
|\mathsf F_\lambda\mathsf F_\eta|\le2.
\]
Direct differentiation gives
\[
\begin{aligned}
\mathsf c_x'(a')
 &=1+
 (1-\mathsf F_\lambda(x)^2)
 \frac{(a')^2(3-(a')^2)}{(1-(a')^2)^2},\\
\mathsf c_x''(a')
 &=(1-\mathsf F_\lambda(x)^2)
 \frac{2a'(3+(a')^2)}{(1-(a')^2)^3}.
\end{aligned}
\]
On \(0\le a'\le1/16\),
\[
1-(a')^2\ge\frac{15}{16},
\qquad
|1-\mathsf F_\lambda^2|\le1,
\]
and hence
\[
|\mathsf c_x|\le\frac1{15},
\qquad
|\mathsf c_x'|
 \le1+\frac{3/256}{(15/16)^2}\le2,
\qquad
|\mathsf c_x''|
 \le\frac{1/2}{(15/16)^3}\le1.
\]
In this notation the interaction coordinate is
\[
\zeta(x)
 =u'\left\{
 a'\mathsf F_\lambda(x)\mathsf F_\eta(x)
 -\nu\kappa\widetilde g_\sigma(x)\mathsf c_x(a')
 \right\}.
\]
The estimates needed for its derivatives are
\[
\begin{aligned}
|\mathsf F_\lambda\mathsf F_\eta
 -\nu\kappa\widetilde g_\sigma\mathsf c_x'|
 &\le 4+\frac18=\frac{33}{8},\\
|\mathsf F_\lambda\mathsf F_\eta
 -\nu\kappa\widetilde g_\sigma\mathsf c_x'|
 &\le 2+\frac18=\frac{17}{8},\\
|a'\mathsf F_\lambda\mathsf F_\eta
 -\nu\kappa\widetilde g_\sigma\mathsf c_x|
 &\le\frac18+\frac1{240}=\frac{31}{240}.
\end{aligned}
\]
The first, coarser bound is sufficient for the first propensity derivative. For every record factor on the amplitude rectangle,
\[
\begin{aligned}
|\partial_{a'}\mathsf r_{\nu,i}|
 &\le2+\frac1{16}\frac{33}{8}\le4,\\
|\partial_{a'}^2\mathsf r_{\nu,i}|
 &\le\frac1{256}\le4,\\
|\partial_{u'}\mathsf r_{\nu,i}|
 &\le2+\frac{31}{240}\le4,\\
|\partial_{a'}\partial_{u'}\mathsf r_{\nu,i}|
 &\le\frac{17}{8}\le4,\\
|\partial_{a'}^2\partial_{u'}\mathsf r_{\nu,i}|
 &\le\frac1{16}\le4,
\qquad
\partial_{u'}^2\mathsf r_{\nu,i}=0.
\end{aligned}
\]
Also \(|\mathsf r_{\nu,i}|\le2\) throughout this closed rectangle.

Assigning the four labelled differentiations to record factors gives at most
\(\mathsf n_{\mathsf c}^4\) terms in the fourth mixed product derivative. Each differentiated factor is bounded by four, and every remaining factor by two. Finite averaging therefore gives
\[
\left|
\partial_{a'}^2\partial_{u'}^2
 \mathsf d_{\mathsf c,\nu,s_{\mathrm{coarse}}}
\right|
 \le256\,\mathsf n_{\mathsf c}^4\,2^{\mathsf n_{\mathsf c}},
\]
and the two-differentiation product rule similarly gives
\[
\left|
\partial_{a'}\partial_{u'}
 \mathsf d_{\mathsf c,\nu,s_{\mathrm{coarse}}}
\right|
 \le16\,\mathsf n_{\mathsf c}^2\,2^{\mathsf n_{\mathsf c}}.
\]
These statements include the empty product, whose derivatives vanish. Taking the even and odd contrasts yields
\[
\left|
\partial_{a'}^2\partial_{u'}^2
 \Delta_{\mathsf c,\mathrm e}
\right|
 \le512\,\mathsf n_{\mathsf c}^4\,2^{\mathsf n_{\mathsf c}},
\qquad
\left|
\partial_{a'}\partial_{u'}
 \Delta_{\mathsf c,\mathrm o}
\right|
 \le16\,\mathsf n_{\mathsf c}^2\,2^{\mathsf n_{\mathsf c}}.
\]

The axis identities and first-order cancellations imply
\[
\partial_{a'}^2\Delta_{\mathsf c,\mathrm e}(a',0)
 =
\partial_{a'}^2\partial_{u'}
 \Delta_{\mathsf c,\mathrm e}(a',0)=0.
\]
Applying the second-order remainder formula first in outcome amplitude and then in propensity amplitude gives
\[
\begin{aligned}
|\Delta_{\mathsf c,\mathrm e}(a,u)|
&\le
512\,\mathsf n_{\mathsf c}^4\,2^{\mathsf n_{\mathsf c}}
 \int_0^a(a-a')\,da'
 \int_0^u(u-u')\,du'\\
&=
128a^2u^2\mathsf n_{\mathsf c}^4\,2^{\mathsf n_{\mathsf c}}.
\end{aligned}
\]
For the odd discrepancy,
\(\partial_{a'}\Delta_{\mathsf c,\mathrm o}(a',0)=0\).
The mixed derivative bound and the mean-value inequality give, successively,
\[
|\partial_{a'}\Delta_{\mathsf c,\mathrm o}(a',u)|
 \le16u\,\mathsf n_{\mathsf c}^2\,2^{\mathsf n_{\mathsf c}},
\qquad
|\Delta_{\mathsf c,\mathrm o}(a,u)|
 \le16au\,\mathsf n_{\mathsf c}^2\,2^{\mathsf n_{\mathsf c}}.
\]

It remains to check singleton cancellation under the actual disclosure. Distinct boundary nodes are separated by \(K/M\ge16\) fine intervals, so at most one active frame node at any covariate is disclosed. Distinct undisclosed coefficient pairs are independent, with each individual coefficient having mean zero. Thus, for every fixed coarse-sign vector,
\[
\mathbb E_{1,\sigma,\delta}[\xi(x)^2]=a^2,
\]
\[
\mathbb E_{1,\sigma,\delta}[\xi(x)\upsilon(x)]
 -\mathbb E_{0,\sigma,\delta}[\xi(x)\upsilon(x)]
 =au\kappa\widetilde g_\sigma(x).
\]
For the first identity, the diagonal terms sum to \(a^2\sum_jf_j(x)^2=a^2\), and every cross term has an undisclosed mean-zero factor. For the second, the undisclosed same-node correlation is
\(\kappa g_\sigma(j/K)\); disclosed same-node terms are common, and
\(g_\sigma(j/K)=0\) at disclosed nodes. Cross terms again vanish. The first moments of \(\xi\) and \(\upsilon\) are common to both branches, and
\[
au\kappa\widetilde g_\sigma(x)
 +t_1(x)(1-a^2)=0.
\]
Consequently every singleton label density agrees with its null density, for either value of its coarse sign. These identities also hold at endpoints by the frame formulas. For the empty subset, normalization of the coefficient weights gives all component densities equal to one. Hence
\[
\Delta_{\mathsf c,\mathrm e}
 =\Delta_{\mathsf c,\mathrm o}=0
\quad\text{if }\mathsf n_{\mathsf c}\le1.
\]

Squaring the discrepancy bounds and using the denominator lower bounds now gives, for \(\mu\)-almost every augmentation and every component,
\[
\mathfrak e_{\mathsf c}
 \le
 2^{14}a^4u^4\mathsf n_{\mathsf c}^{\,8}
       8^{\mathsf n_{\mathsf c}}
 \mathbf1\{\mathsf n_{\mathsf c}\ge2,\
           \mathsf k_{\mathsf c}\ge2\},
\]
\[
\mathfrak o_{\mathsf c}
 \le
 2^8a^2u^2\mathsf n_{\mathsf c}^{\,4}
       8^{\mathsf n_{\mathsf c}}
 \mathbf1\{\mathsf n_{\mathsf c}\ge2,\
           \mathsf k_{\mathsf c}\ge1\}.
\]
% lean: CausalSmith.Stat.FinitepHomogeneityDensegamma.component_discrepancy_bounds
% lean: CausalSmith.Stat.FinitepHomogeneityDensegamma.component_activity_envelopes

\item
We bound the expectations of the budgets. Define the two occupancy scores by
\[
\mathsf S_{\mathrm{one}}(\omega)
 =
 \sum_{\mathsf c\in\mathscr C(\omega)}
 \mathsf n_{\mathsf c}^{\,8}8^{\mathsf n_{\mathsf c}}
 \mathbf1\{\mathsf n_{\mathsf c}\ge2,\
           \mathsf k_{\mathsf c}\ge2\},
\]
\[
\begin{aligned}
\mathsf S_{\mathrm{pair}}(\omega)
 =\frac12
 \sum_{\mathsf c,\mathsf c'\in\mathscr C(\omega)}
 &\mathsf n_{\mathsf c}^{\,4}\mathsf n_{\mathsf c'}^{\,4}
   8^{\mathsf n_{\mathsf c}+\mathsf n_{\mathsf c'}}\\
 &\times
 \mathbf1\left\{
 \begin{array}{c}
 \mathsf c\ne\mathsf c',\
 \mathsf j(\mathsf c)=\mathsf j(\mathsf c'),\\
 \mathsf n_{\mathsf c},\mathsf n_{\mathsf c'}\ge2,\
 \mathsf k_{\mathsf c},\mathsf k_{\mathsf c'}\ge1
 \end{array}\right\}.
\end{aligned}
\]
The component estimates imply
\[
\mathcal A\le2^{14}a^4u^4\mathsf S_{\mathrm{one}},
\qquad
\mathcal B\le2^{16}a^4u^4\mathsf S_{\mathrm{pair}}
\quad\mu\text{-almost surely}.
\]

Under \(\mu\), the fine-cell indices are independent uniform draws from
\(\{0,\ldots,K-1\}\), and the flags are independent Bernoulli draws, independent of those indices. A connected component of size \(r_{\mathrm{occ}}\) occupies consecutive fine cells spanning at most \(r_{\mathrm{occ}}\) cells. For a specified observation subset of that size and a specified starting cell, the necessary occupancy event has probability at most
\[
\left(\frac{r_{\mathrm{occ}}}{K}\right)^{r_{\mathrm{occ}}}.
\]
A window extending past the final cell has fewer available cells and obeys the same bound. The probability of at least two marks in that subset is at most
\(r_{\mathrm{occ}}^2\varepsilon^2\). There are at most \(K\) starting cells and
\(\binom n{r_{\mathrm{occ}}}\) observation subsets. Overcounting these necessary events yields
\[
\int\mathsf S_{\mathrm{one}}\,d\mu
 \le
 \varepsilon^2K
 \sum_{r_{\mathrm{occ}}=2}^{n}
 \binom n{r_{\mathrm{occ}}}
 r_{\mathrm{occ}}^{10}8^{r_{\mathrm{occ}}}
 \left(\frac{r_{\mathrm{occ}}}{K}\right)^{r_{\mathrm{occ}}}.
\]

For two distinct components, their observation subsets are disjoint. Each requires a mark, so for sizes
\(r_{\mathrm{occ}},r'_{\mathrm{occ}}\) the joint flag probability is at most
\(r_{\mathrm{occ}}r'_{\mathrm{occ}}\varepsilon^2\). A coarse pair contains
\(2K/M\) fine cells. Thus the number of ordered starting-cell choices in a common coarse pair is at most
\[
\frac M2\left(\frac{2K}{M}\right)^2=\frac{2K^2}{M}.
\]
The number of disjoint subset choices is bounded by
\(\binom n{r_{\mathrm{occ}}}\binom n{r'_{\mathrm{occ}}}\).
Using independence on the disjoint subsets, and overcounting by ordered pairs, gives
\[
\int\mathsf S_{\mathrm{pair}}\,d\mu
 \le
 \frac{2\varepsilon^2K^2}{M}
 \left\{
 \sum_{r_{\mathrm{occ}}=2}^{n}
 \binom n{r_{\mathrm{occ}}}
 r_{\mathrm{occ}}^{5}8^{r_{\mathrm{occ}}}
 \left(\frac{r_{\mathrm{occ}}}{K}\right)^{r_{\mathrm{occ}}}
 \right\}^{2}.
\]
When the collection of eligible pairs is empty, its score is zero and this bound still applies.

Define the occupancy ratio by
\[
\zeta_{\mathrm{occ}}=8e\,\frac nK.
\]
The factorial and binomial bounds
\[
r_{\mathrm{occ}}!\ge
 \left(\frac{r_{\mathrm{occ}}}{e}\right)^{r_{\mathrm{occ}}},
\qquad
\binom n{r_{\mathrm{occ}}}
 \le\frac{n^{r_{\mathrm{occ}}}}{r_{\mathrm{occ}}!}
\]
imply, for every \(r_{\mathrm{occ}}\ge2\) and every nonnegative integer
\(d_{\mathrm{occ}}\),
\[
\binom n{r_{\mathrm{occ}}}
 r_{\mathrm{occ}}^{d_{\mathrm{occ}}}8^{r_{\mathrm{occ}}}
 \left(\frac{r_{\mathrm{occ}}}{K}\right)^{r_{\mathrm{occ}}}
 \le
 r_{\mathrm{occ}}^{d_{\mathrm{occ}}}
 \zeta_{\mathrm{occ}}^{r_{\mathrm{occ}}}
 \le
 r_{\mathrm{occ}}^{d_{\mathrm{occ}}+1}
 \zeta_{\mathrm{occ}}^{r_{\mathrm{occ}}}.
\]
The factorial inequality follows from the usual Stirling lower bound, whose square-root factor is at least one for positive integers. Extending the nonnegative finite sums gives
\[
\int\mathcal A\,d\mu
 \le
 2^{14}a^4u^4\varepsilon^2K
 \sum_{r_{\mathrm{occ}}=2}^{\infty}
 r_{\mathrm{occ}}^{11}
 \zeta_{\mathrm{occ}}^{r_{\mathrm{occ}}},
\]
\[
\int\mathcal B\,d\mu
 \le
 \frac{2^{17}a^4u^4\varepsilon^2K^2}{M}
 \left\{
 \sum_{r_{\mathrm{occ}}=2}^{\infty}
 r_{\mathrm{occ}}^{6}
 \zeta_{\mathrm{occ}}^{r_{\mathrm{occ}}}
 \right\}^{2}.
\]

For \(d_{\mathrm{ser}}\in\{6,11\}\), the elementary inequality
\(r_{\mathrm{occ}}\le2^{r_{\mathrm{occ}}-1}\) gives
\[
\begin{aligned}
\sum_{r_{\mathrm{occ}}=2}^{\infty}
 r_{\mathrm{occ}}^{d_{\mathrm{ser}}}
 \zeta_{\mathrm{occ}}^{r_{\mathrm{occ}}}
 &\le
 2^{d_{\mathrm{ser}}}\zeta_{\mathrm{occ}}^2
 \sum_{r_{\mathrm{occ}}=2}^{\infty}
 (2^{d_{\mathrm{ser}}}\zeta_{\mathrm{occ}})^{r_{\mathrm{occ}}-2}\\
 &=
 \frac{2^{d_{\mathrm{ser}}}\zeta_{\mathrm{occ}}^2}
      {1-2^{d_{\mathrm{ser}}}\zeta_{\mathrm{occ}}}
 \le2^{d_{\mathrm{ser}}+1}\zeta_{\mathrm{occ}}^2.
\end{aligned}
\]
The geometric series converges because
\[
0\le\zeta_{\mathrm{occ}}\le24\,\frac nK,
\qquad
2^{11}\zeta_{\mathrm{occ}}
 \le2^{11}\cdot24\cdot2^{-40}<\frac12.
\]
Here we used \(e<3\) and \(2^{40}n\le K\). Consequently,
\[
\begin{aligned}
\int\mathcal A\,d\mu
 &\le
 2^{26}24^2\,
 \frac{a^4u^4\varepsilon^2n^2}{K}
 \le
 \frac{2^{38}a^4u^4\varepsilon^2n^2}{K},\\
\int\mathcal B\,d\mu
 &\le
 2^{31}24^4\,
 \frac{a^4u^4\varepsilon^2n^4}{MK^2}
 \le
 \frac{2^{56}a^4u^4\varepsilon^2n^4}{MK^2},
\end{aligned}
\]
since \(2^{26}24^2<2^{38}\) and \(2^{31}24^4<2^{56}\).
The budgets are measurable finite sums. The uniform component denominator bounds and finite label alphabet make them bounded for fixed parameters, hence integrable under the probability law \(\mu\).
% lean: CausalSmith.Stat.FinitepHomogeneityDensegamma.component_occupancy_budgets

\item\label{proof:full-record-copula-component-bound:record-projections}
We identify the original-record projections. Define the recovery map by
\[
\mathsf O_L(\omega,\mathbf z)
 =
 \left(
 X_i,\frac{1+z_i^{\mathrm t}}2,B_i z_i^{\mathrm y}L
 \right)_{i=1}^{n}.
\]
This map is measurable.

For a fixed \(\sigma\), the coefficient law in
\cref{obj:def:copula-frame-priors} assigns mass
\[
\prod_{j=0}^{K}
 \frac{1+\nu\kappa g_\sigma(j/K)\lambda_j\eta_j}{4}
\]
to a coefficient draw. At boundary nodes, \(g_\sigma(j/K)=0\). Its disclosure law is therefore \(\mu_{\partial}\), and conditioning on a disclosure gives precisely the weights
\(\mathsf w_{\nu,\sigma,\delta}\). In particular, for every function \(\Psi\) on the coefficient space,
\[
\int
 \mathbb E_{\nu,\sigma,\delta}[\Psi]\,
 \mu_{\partial}(d\delta)
 =
 \sum_{\boldsymbol\lambda,\boldsymbol\eta}
 \left\{
 \prod_{j=0}^{K}
 \frac{1+\nu\kappa g_\sigma(j/K)\lambda_j\eta_j}{4}
 \right\}
 \Psi(\boldsymbol\lambda,\boldsymbol\eta).
\]
For the null branch, the component factorization in \cref{proof:full-record-copula-component-bound:conditional-comparisons} identifies
\(\mathsf f_0\) with the conditional full-record label average. For the alternative branch, that average is the definition of \(\mathsf f_1\). Integrating the disclosure thus recovers the respective original coefficient averages.

For a fixed latent draw and covariate \(x\), the flag and fair-label weights give
\[
\Pr_\nu\left(
 A=\frac{1+z^{\mathrm t}}2,\,
 Y=z^{\mathrm y}L\mid x
 \right)
 =
 \frac{\varepsilon}{4}
 \left\{
 1+z^{\mathrm t}\xi(x)+z^{\mathrm y}\upsilon(x)
   +z^{\mathrm t}z^{\mathrm y}\zeta(x)
 \right\},
\]
and, summing the retained outcome-sign label when the flag is zero,
\[
\Pr_\nu\left(
 A=\frac{1+z^{\mathrm t}}2,\,
 Y=0\mid x
 \right)
 =
 \sum_{z^{\mathrm y}\in\{-1,1\}}
 \frac{1-\varepsilon}{4}(1+z^{\mathrm t}\xi(x))
 =
 \frac{1-\varepsilon}{2}(1+z^{\mathrm t}\xi(x)).
\]
These are exactly the categories of \cref{obj:def:marked-table}. Independent record generation conditional on the latent draw, followed by its prior average, therefore proves
\[
(\mu\otimes P_0)\circ\mathsf O_L^{-1}=\mathbb P_0,
\qquad
(\mu\otimes P_1)\circ\mathsf O_L^{-1}=\mathbb P_1.
\]
% lean: CausalSmith.Stat.FinitepHomogeneityDensegamma.null_augmented_projection_label_bind
% lean: CausalSmith.Stat.FinitepHomogeneityDensegamma.alternative_augmented_projection_label_bind
% lean: CausalSmith.Stat.FinitepHomogeneityDensegamma.prior_label_design_bind_eq_mixture

\item
Assume the sum of the two stated public budget bounds is at most \(2^{-16}\). We first obtain an envelope for conditional total variation.

On a finite alphabet, absolute continuity and Cauchy--Schwarz give
\[
\operatorname{TV}(Q(\omega),P_0(\omega))
 \le\frac12\sqrt{\chi^2(Q(\omega)\Vert P_0(\omega))},
\]
\[
\operatorname{TV}(P_1(\omega),Q(\omega))
 \le\frac12\sqrt{\chi^2(P_1(\omega)\Vert Q(\omega))}.
\]
Indeed, total variation is half the sum of absolute differences of the masses; writing those differences relative to the dominating law and applying Cauchy--Schwarz gives the displayed bounds.

For \(0\le x_{\mathrm{exp}}\le1/256\), comparison of the exponential series with the geometric series gives
\[
e^{x_{\mathrm{exp}}}
 \le\frac1{1-x_{\mathrm{exp}}}
 \le1+2x_{\mathrm{exp}},
\]
where the second inequality follows after multiplication by
\(1-x_{\mathrm{exp}}>0\). If
\(\mathcal A(\omega)+\mathcal B(\omega)\le1/256\), the conditional divergence bounds and the triangle inequality therefore yield
\[
\begin{aligned}
\operatorname{TV}(P_0(\omega),P_1(\omega))
 &\le
 \frac12\sqrt{2\mathcal A(\omega)}
 +\frac12\sqrt{2\mathcal B(\omega)}\\
 &\le\sqrt{\mathcal A(\omega)+\mathcal B(\omega)}
 \le\frac1{16}.
\end{aligned}
\]
The middle inequality follows by squaring and using
\(2\sqrt{\mathcal A\mathcal B}\le\mathcal A+\mathcal B\).
In the complementary case,
\(\mathcal A(\omega)+\mathcal B(\omega)>1/256\), total variation between probability laws is at most one. Both cases consequently give
\[
\operatorname{TV}(P_0(\omega),P_1(\omega))
 \le\frac1{16}
      +256\bigl(\mathcal A(\omega)+\mathcal B(\omega)\bigr)
\quad\mu\text{-almost surely}.
\]

For any measurable augmented event
\(\mathsf E_{\mathrm{aug}}\subseteq
 \mathsf\Omega_{\mathrm{aug}}\times\mathsf Z_n\),
define its label section by
\[
\mathsf E_{\mathrm{aug},\omega}
 =\{\mathbf z:(\omega,\mathbf z)\in\mathsf E_{\mathrm{aug}}\}.
\]
The common augmentation marginal implies
\[
\begin{aligned}
&|(\mu\otimes P_0)(\mathsf E_{\mathrm{aug}})
  -(\mu\otimes P_1)(\mathsf E_{\mathrm{aug}})|\\
&\qquad\le
 \int
 |P_0(\omega)(\mathsf E_{\mathrm{aug},\omega})
  -P_1(\omega)(\mathsf E_{\mathrm{aug},\omega})|\,\mu(d\omega)\\
&\qquad\le
 \frac1{16}
 +256\left(\int\mathcal A\,d\mu+\int\mathcal B\,d\mu\right)
 \le\frac1{16}+256\cdot2^{-16}
 =\frac{17}{256}.
\end{aligned}
\]
Taking the supremum over augmented events proves
\[
\operatorname{TV}(\mu\otimes P_0,\mu\otimes P_1)
 \le\frac{17}{256}.
\]
% lean: CausalSmith.Stat.FinitepHomogeneityDensegamma.copula_augmented_tv_budget

For every measurable original-record event
\(\mathsf E_{\mathrm{rec}}\), its inverse image under \(\mathsf O_L\) is an augmented event. The projection identities from \cref{proof:full-record-copula-component-bound:record-projections} therefore give
\[
\operatorname{TV}(\mathbb P_0,\mathbb P_1)
 \le
 \operatorname{TV}(\mu\otimes P_0,\mu\otimes P_1)
 \le\frac{17}{256}<\frac18.
\]
% lean: CausalSmith.Stat.FinitepHomogeneityDensegamma.copula_tv_map_le

\item
Finally, we verify that the augmentation is induced by the stipulated latent construction. Define its latent coefficient law by the masses
\[
\mathsf W_\nu
 \bigl(\{\sigma,\boldsymbol\lambda,\boldsymbol\eta\}\bigr)
 =
 2^{-M/2}
 \prod_{j=0}^{K}
 \frac{1+\nu\kappa g_\sigma(j/K)\lambda_j\eta_j}{4}.
\]
Define the joint latent augmented law by
\[
\begin{aligned}
&\mathsf L_\nu
 (d\sigma,d\boldsymbol\lambda,d\boldsymbol\eta,
  d\mathbf X,d\mathbf B,d\delta,d\mathbf z)\\
&\quad=
 \mathsf W_\nu(d\sigma,d\boldsymbol\lambda,d\boldsymbol\eta)\,
 \lambda^{\otimes n}(d\mathbf X)\,
 \operatorname{Bernoulli}(\varepsilon)^{\otimes n}(d\mathbf B)\,
 \operatorname{Dirac}_{\mathsf D(\boldsymbol\lambda,\boldsymbol\eta)}
       (d\delta)\,
 \left(\prod_{i=1}^{n}\mathsf r_{\nu,i}\right)
 \mathsf F_n(d\mathbf z).
\end{aligned}
\]
The coefficient weights sum to one, and the conditional label product integrates to one. Thus
\[
\mathsf L_\nu(\text{entire space})=1,
\qquad
\mathsf L_\nu\{\delta=\mathsf D(\boldsymbol\lambda,\boldsymbol\eta)\}=1.
\]
Conditioning its finite coefficient average on the disclosure, and then using the component factorization for branch zero and the full coefficient average for branch one, gives
\[
\mathsf L_\nu\circ
 \bigl((\sigma,\boldsymbol\lambda,\boldsymbol\eta,\omega,\mathbf z)
       \mapsto(\omega,\mathbf z)\bigr)^{-1}
 =\mu\otimes P_\nu,
\qquad \nu\in\{0,1\}.
\]

There are \(M+1\) disclosed nodes, and the coarse tents vanish at all of them. Their coefficient pairs are therefore independent fair pairs conditional on every coarse-sign vector. Summing the undisclosed coefficient weights gives, for each
\(\sigma_*\in\{-1,1\}^{M/2}\) and each actual boundary disclosure \(\delta\),
\[
\Pr_\nu\{\sigma=\sigma_*,\mathcal D_{\partial}=\delta\}
 =2^{-M/2}4^{-(M+1)}.
\]
Define the fair coarse-sign law by
\[
\mathsf U_\sigma
 =2^{-M/2}\sum_{\sigma\in\{-1,1\}^{M/2}}
       \operatorname{Dirac}_{\sigma}.
\]
The covariates and flags are drawn independently of these latent variables, and integrating the normalized labels leaves their law unchanged. Hence
\[
\mathsf L_\nu\circ
 \bigl((\sigma,\boldsymbol\lambda,\boldsymbol\eta,\omega,\mathbf z)
       \mapsto(\sigma,\omega)\bigr)^{-1}
 =\mathsf U_\sigma\otimes\mu.
\]
This proves both the independent fair-sign assertion and its product relation with the full augmentation. Since \(L>0\), recovery also satisfies
\(B_i=\mathbf1\{Y_i\ne0\}\). Together with the kernel, budget, projection, and total-variation properties already established, these are all the required conclusions.
% lean: CausalSmith.Stat.FinitepHomogeneityDensegamma.copulaLatentAugmentedLaw_map_conditional
% lean: CausalSmith.Stat.FinitepHomogeneityDensegamma.copulaLatentAugmentedLaw_sign_independent
% lean: CausalSmith.Stat.FinitepHomogeneityDensegamma.full_record_copula_component_bound
\end{enumerate}
\end{proof}
